% =================================================================================================
% File:			nome_del_documento.tex
% Description:	Defiinisce la struttura del documento di Nome del documento 
% Created:		2014/12/05
% Author:		Santacatterina Luca
% Email:		s88.luca@gmail.com
% =================================================================================================
% Modification History:
% Version		Modifier Date		Change											Author
% 0.0.1 		2014/12/05 			stesura documento di prova						Luca S.
% =================================================================================================
%

% TEMPLATE DA USARE AD OGNI NUOVO DOCUMENTO

% =================================================================================================
% File:			package.tex
% Description:	Defiinisce i contenuti obbligatori della prima pagina
% Created:		2014/12/05
% Author:		Santacatterina Luca
% Email:		s88.luca@gmail.com
% =================================================================================================
% Modification History:
% Version		Modifier Date		Change											Author
% 0.0.1 		2014/12/05 			Prima emissione 								Luca S.
% =================================================================================================
%

\documentclass[11pt,a4paper]{article}
\usepackage[left=3.5cm,right=3.5cm,top=3.0cm,bottom=2.5cm]{geometry}
\usepackage[font=small,labelfont=bf]{caption}
\usepackage[italian]{babel}
\usepackage[utf8]{inputenc}
\usepackage[T1]{fontenc}

\usepackage{color}
\usepackage{enumitem}
\usepackage{eurosym}
\usepackage{fancyhdr}
\usepackage{geometry}
\usepackage{graphicx}
\usepackage{hyperref}
\usepackage{lastpage}
\usepackage{lastpage}
\usepackage{listings}
\usepackage{longtable}
\usepackage{multirow}
\usepackage{pdflscape}
\usepackage{pdfpages}
\usepackage{subcaption}
\usepackage{tabularx}
\usepackage{tikz}
\usepackage{titlesec}
\usepackage{verbatim}
\usepackage{xr}

\pagestyle{fancy}
\hypersetup{
	colorlinks=true,
	linkcolor=black,
	urlcolor=black
}



% =================================================================================================
% File:			commands.tex
% Description:	Defiinisce i comandi personalizzati utilizzati per redarre i documenti
% Created:		2014-12-05
% Author:		Santacatterina Luca
% Email:		s88.luca@gmail.com
% =================================================================================================
% Modification History:
% Version		Modifier Date		Change												Author
% 0.0.1 		2014-12-05 			Prima emissione 									Luca S.
% =================================================================================================
% 0.0.2 		2014-12-13 			Inseriti comandi nome e la vers. dei file 			Luca S.
% =================================================================================================
% 0.0.3			2014-12-19			inserito comando per richiamare il libro di testo	Tesser Paolo
% =================================================================================================
% 0.0.4			2015-01-02			inserito comando per le revisioni					Tesser Paolo
% =================================================================================================
% 0.0.5			2015-01-07			inseriti comandi per i documenti mancanti			Tesser Paolo
% =================================================================================================
% 0.0.6			2015-01-12			cambiata versione dei documenti a 1.0.0				Tesser Paolo
% =================================================================================================
% 0.0.7			2015-01-14			comandi per la generazione di grafici a torta/barre	Tesser Paolo

% DEFINIZIONE COMANDI BASE CHE POI ANDRANNO RIDEFINITI

% ==================================================================================================
% COMANDI DA RIDEFINIRE
% ==================================================================================================

% Nome/Versione/Data del documento
\newcommand{\documentName}{ERRORE}
\newcommand{\documentVersion}{ERRORE}
\newcommand{\documentDate}{ERRORE}

% Editori del documento
\newcommand{\documentEditors}{ERRORE}

% Verificatore del documento
\newcommand{\documentVerifiers}{ERRORE}

% Approvazione del documento
\newcommand{\documentApprovers}{ERRORE}

% Lista di distribuzione del documento
\newcommand{\documentDistributionList}{\groupName\\\commitNameM\\\commitNameS\\\proposerName}

% Uso del documento
\newcommand{\documentUsage}{ERRORE}

% Sommario del documento
\newcommand{\documentSummary}{ERRORE}

% Vesione dei documenti
\newcommand{\docVersionAdR}{\emph{v1.0.0}} % Analisi dei Requisiti
\newcommand{\docVersionGlo}{\emph{v1.0.0}} % Glossario
\newcommand{\docVersionNdP}{\emph{v1.0.0}} % Norme di Progetto
\newcommand{\docVersionPdP}{\emph{v1.0.0}} % Piano di Progetto
\newcommand{\docVersionPdQ}{\emph{v1.0.0}} % Piano di Qualifica
\newcommand{\docVersionSdF}{\emph{v1.0.0}} % Studio di Fattibilità
\newcommand{\docVersionST}{\emph{v1.0.0}} % Specifica Tecnica
\newcommand{\docVersionDdP}{\emph{v1.0.0}} % Definizione di Prodotto
\newcommand{\docVersionMU}{\emph{v1.0.0}} % Manuale Utente

% ==================================================================================================
% COMANDI DA NON RIDEFINIRE
% ==================================================================================================

% Nome del progetto
\newcommand{\projectName}{BDSMApp}
\newcommand{\groupEmail}{\textit{\href{mailto:info@mashup-unipd.it}{info@mashup-unipd.it}}}
\newcommand{\groupName}{\emph{MashUp}}

% Ruoli del progetto
\newcommand{\roleProjectManager}{\emph{Responsabile di Progetto}}
\newcommand{\roleAdministrator}{\emph{Amministratore di Progetto}}
\newcommand{\roleAnalyst}{\emph{Analista}}
\newcommand{\roleProgrammer}{\emph{Programmatore}}
\newcommand{\roleDesigner}{\emph{Progettista}}
\newcommand{\roleVerifier}{\emph{Verificatore}}

% Referenti e commitenti (M = master, S = slave)
\newcommand{\proposerName}{\emph{Dott. David Santucci} - Zing Srl}
\newcommand{\commitNameM}{\emph{Prof. Tullio Vardanega}}
\newcommand{\commitNameS}{\emph{Prof. Riccardo Cardin}}

% Abbreviazioni per richiamare il nome esteso dei diversi documenti da redarre
\newcommand{\docGlossary}{\emph{Glossario}} 
\newcommand{\docAdR}{\emph{Analisi dei Requisiti \currentVersion}}
\newcommand{\docNdP}{\emph{Norme di Progetto \currentVersion}}
\newcommand{\docPdP}{\emph{Piano di Progetto \currentVersion}}
\newcommand{\docPdQ}{\emph{Piano di Qualifica \currentVersion}}
\newcommand{\docSdF}{\emph{Studio di Fattibilità \currentVersion}}

% Nome dei documenti
\newcommand{\docNameAdR}{\emph{Analisi dei Requisiti}} % Analisi dei Requisiti
\newcommand{\docNameGlo}{\emph{Glossario}} % Glossario
\newcommand{\docNameNdP}{\emph{Norme di Progetto}} % Norme di Progetto
\newcommand{\docNamePdP}{\emph{Piano di Progetto}} % Piano di Progetto
\newcommand{\docNamePdQ}{\emph{Piano di Qualifica}} % Piano di Qualifica
\newcommand{\docNameSdF}{\emph{Studio di Fattibilità}} % Studio di Fattibilità
\newcommand{\docNameST}{\emph{Specifica Tecnica}} % Specifica Tecnica
\newcommand{\docNameDdP}{\emph{Definizione di Prodotto}} % Definizione di Prodotto
\newcommand{\docNameMU}{\emph{Manuale Utente}} % Manuale Utente

% Nome e versione dei documenti
\newcommand{\docNameVersionAdR}{\docNameAdR{} \docVersionAdR} % Analisi dei Requisiti
\newcommand{\docNameVersionGlo}{\docNameGlo{} \docVersionGlo} % Glossario
\newcommand{\docNameVersionNdP}{\docNameNdP{} \docVersionNdP} % Norme di Progetto
\newcommand{\docNameVersionPdP}{\docNamePdP{} \docVersionPdP} % Piano di Progetto
\newcommand{\docNameVersionPdQ}{\docNamePdQ{} \docVersionPdQ} % Piano di Qualifica
\newcommand{\docNameVersionSdF}{\docNameSdF{} \docVersionSdF} % Studio di Fattibilità
\newcommand{\docNameVersionST}{\docNameST{} \docVersionST} % Specifica Tecnica
\newcommand{\docNameVersionDdP}{\docNameDdP{} \docVersionDdP} % Definizione di Prodotto
\newcommand{\docNameVersionMU}{\docNameMU{} \docVersionMU} % Manuale Utente

% Nome delle revisioni
\newcommand{\RR}{\emph{Revisione dei Requisiti}}
\newcommand{\RPmin}{\emph{Revisione di Progettazione minima}}
\newcommand{\RPmax}{\emph{Revisione di Progettazione massima}}
\newcommand{\RQ}{\emph{Revisione di Qualifica}}
\newcommand{\RA}{\emph{Revisione di Accettazione}}

% Descrizione scopo del glossario
\newcommand{\glossarioDesc}{Al fine di evitare ogni ambiguità relativa al linguaggio usato nei documenti viene allegato il \docNameVersionGlo.\\
Esso ha lo scopo di definire ed analizzare tutti i termini tecnici, di progetto e di ambiguità fornendo una ben precisa descrizione.\\
Tutte le occorrenze presenti nei documenti verranno contrassegnate con una ``G'' in formato pedice.}

% Descrizione dello scopo del prodtto
\newcommand{\productScope}{Lo scopo del prodotto è di creare una nuova infrastruttura che permetta di interrogare Big Data recuperati dai social network, quali: Facebook, Twitter, Instagram.
L'applicazione sarà composta da due parti:
	\begin{itemize}
		\item consultazione e interrogazione con interfaccia web per utente;
		\item servizi web REST interrogabili.
	\end{itemize}
}

\newcommand{\sommerville}{Software Engineering - Ian Sommerville - 9th Edition(2010)}

% Comandi per la generazione di grafici a torta e a barre
\newcommand{\pie}[3][]{
    \begin{scope}[#1]
    \pgfmathsetmacro{\curA}{90}
    \pgfmathsetmacro{\r}{1}
    \def\c{(0,0)}
    \node[pie title] at (90:1.3) {#2};
    \foreach \v\s in{#3}{
        \pgfmathsetmacro{\deltaA}{\v/100*360}
        \pgfmathsetmacro{\nextA}{\curA + \deltaA}
        \pgfmathsetmacro{\midA}{(\curA+\nextA)/2}

        \path[slice,\s] \c
            -- +(\curA:\r)
            arc (\curA:\nextA:\r)
            -- cycle;
        \pgfmathsetmacro{\d}{max((\deltaA * -(.5/50) + 1) , .5)}

        \begin{pgfonlayer}{foreground}
        \path \c -- node[pos=\d,pie values,values of \s]{$\v\%$} +(\midA:\r);
        \end{pgfonlayer}

        \global\let\curA\nextA
    }
    \end{scope}
}

\newcommand{\legend}[2][]{
    \begin{scope}[#1]
    \path
        \foreach \n/\s in {#2}
            {
                  ++(0,-5pt) node[\s,legend box] {} +(5pt,0) node[legend label] {\n} % ++(x,y) -> y: spazio tra i vari campi della legenda, x: l'obliquità
            }
    ;
    \end{scope}
}

% Comandi per il glossario
\newcommand{\gloss}{\small\ped{\textbf{G}}}
% =================================================================================================
% File:			package.tex
% Description:	Defiinisce  lo stile che viene utilizzato nelle pagine dei documenti
% Created:		2014/12/05
% Author:		Santacatterina Luca
% Email:		s88.luca@gmail.com
% =================================================================================================
% Modification History:
% Version		Modifier Date		Change											Author
% 0.0.1 		2014/12/05 			Prima emissione 								Luca S.
% =================================================================================================
%

\fancypagestyle{styleFirstPage}{

}


\fancypagestyle{styleIndexPage}{	
	\pagestyle{empty}
	% HEADER
	\fancyhead{}
	\setlength{\topmargin}{-65pt}
	\setlength{\headsep}{65pt}
	\fancyhead[L]{\definecolor{titlepagecolor}{cmyk}{0,0,0,0.40}
	\begin{tikzpicture}[remember picture,overlay]
		\path [fill=titlepagecolor](-20,1.8) rectangle(1\paperwidth,-0.3);
	\end{tikzpicture}
	\includegraphics[width=150pt]{../template/images/logo/logo_orizzontale.pdf}}
	\fancyhead[R]{\uppercase{\leftmark}}
	% FOOTER
	\fancyfoot[L]{}
	\fancyfoot[R]{}
	\fancyfoot[C]{\thepage}
	\pagenumbering{Roman}
	\renewcommand{\footrulewidth}{0.0pt}
	\renewcommand{\headrulewidth}{0.0pt}
}


\fancypagestyle{styleDocPages}{
	\thispagestyle{empty}
	\thispagestyle{fancy}
	\pagestyle{fancy}
	% HEADER
	\fancyhead{}
	\fancyhead[L]{\definecolor{titlepagecolor}{cmyk}{0,0,0,0.40}
	\begin{tikzpicture}[remember picture,overlay]
		\path [fill=titlepagecolor](-20,1.8) rectangle(1\paperwidth,-0.3);
	\end{tikzpicture}
	\includegraphics[width=150pt]{../template/images/logo/logo_orizzontale.pdf}}
	\fancyhead[R]{\uppercase{\leftmark}}
	% FOOTER
	\fancyfoot{}
	\fancyfoot[L]{\documentName}
	\fancyfoot[R]{Pagina \thepage~di~\pageref{LastPage}}
	\setcounter{page}{0}
	\pagenumbering{arabic}
	\renewcommand{\headrulewidth}{0.0pt}
  	\renewcommand{\footrulewidth}{0.2pt}
}

\begin{document}
	% COMANDI DA RIDEFINIRE
<<<<<<< HEAD
	\renewcommand{\documentName}{Verbale del 2014-11-25}
	\renewcommand{\documentDate}{2014-12-02}
	\renewcommand{\documentEditors}{Ceccon Lorenzo}
	\renewcommand{\documentVerifiers}{Faccin Nicola}
	\renewcommand{\documentApprovers}{Paolo Tesser}
	\renewcommand{\documentUsage}{Interno}
	\renewcommand{\documentSummary}{Questo documento riassume il primo incontro tra i membri del gruppo Mashup. }
=======
	\renewcommand{\documentName}{Analisi dei Requisiti}
	\renewcommand{\documentVersion}{v0.0.1}
	\renewcommand{\documentDate}{02/12/2014}
	\renewcommand{\documentEditors}{Cognome Nome}
	\renewcommand{\documentVerifiers}{Cognome Nome\\Cognome Nome}
	\renewcommand{\documentApprovers}{Cognome Nome\\Cognome Nome}
	\renewcommand{\documentUsage}{Esterno}
	\renewcommand{\documentSummary}{Breve descrizione del documento in questione. Serve per dare una idea di massima sul contenuto del documento. L'intera lunghezza della frase non dovrà superare le 3/4 righe totali.}
>>>>>>> develop

	% FRONTESPIZIO
	\pagenumbering{gobble}
	\pagestyle{styleFirstPage}
	% =================================================================================================
% File:			firstpage.tex
% Description:	Defiinisce i contenuti obbligatori della prima pagina
% Created:		2014/12/05
% Author:		Santacatterina Luca
% Email:		s88.luca@gmail.com
% =================================================================================================
% Modification History:
% Version		Modifier Date		Change											Author
% 0.0.1 		2014/12/05 			Prima emissione 								Luca S.
% =================================================================================================
%


% DEFINIZIONE DELLA PRIMA PAGINA

\thispagestyle{empty}
\setlength{\topmargin}{-60pt}
% Logo del gruppo
\begin{figure}[!t]
	\centering
	\includegraphics[width=10cm]{../template/images/logo/zing.pdf}
\end{figure}
\vspace{1pt}
\begin{figure}[!t]
	\centering
	\includegraphics[width=8cm]{../template/images/logo/logo_verticale.pdf}
\end{figure}

% Email del gruppo
\begin{center}
 	\groupEmail
\end{center}

% Informazioni del documento
\begin{center}
\vspace{1cm}
	\textbf{\Large Informazioni Documento} \\[10pt]
	\begin{longtable}{l|p{150pt}}
		\textbf{Nome documento} & \documentName \\ & \\
<<<<<<< HEAD
<<<<<<< HEAD
<<<<<<< HEAD
=======
>>>>>>> a832c2ccc18967166c62b459de65c98f82df83cf
		\textbf{Data redazione} & \documentDate \\ & \\
		\textbf{Redattori} & \parbox[t]{\linewidth}{\documentEditors} \\ & \\
		\textbf{Verificatori} & \parbox{\linewidth}{\documentVerifiers} \\ & \\
		\textbf{Approvazione} & \parbox{\linewidth}{\documentApprovers} \\ & \\
		\textbf{Lista distribuzione} & \parbox{\linewidth}{\documentDistributionList} \\ & \\
		\textbf{Uso} & \parbox{\linewidth}{\documentUsage} \\
<<<<<<< HEAD
=======
=======
>>>>>>> develop
		\textbf{Versione} & \documentVersion \\ & \\
		\textbf{Data redazione} & \documentDate \\ & \\
		\textbf{Redattori} & \parbox[t]{\textwidth}{\documentEditors} \\ & \\
		\textbf{Verificatori} & \parbox{\textwidth}{\documentVerifiers} \\ & \\
		\textbf{Approvazione} & \parbox{\textwidth}{\documentApprovers} \\ & \\
		\textbf{Lista distribuzione} & \parbox{\textwidth}{\documentDistributionList} \\ & \\
		\textbf{Uso} & \parbox{\textwidth}{\documentUsage} \\
<<<<<<< HEAD
>>>>>>> develop
=======
>>>>>>> develop
=======
>>>>>>> a832c2ccc18967166c62b459de65c98f82df83cf
	\end{longtable}
	\vspace{10pt}
	\textbf{\Large Sommario} \\
\end{center}
\documentSummary
	\pagebreak


	% PAGINA REVISIONI & PAGINA INDICI
	\setcounter{page}{1}
	\pagestyle{styleIndexPage}
	% =================================================================================================
% File:			norme_di_progetto.tex
% Description:	Definisce le modifiche apportate al documento
% Created:		2014-12-10
% Author:		Tesser Paolo
% Email:		tesser.paolo@mashup-unipd.it
% =================================================================================================
% Modification History:
% Version		Modifier Date		Change											Author
% 0.0.1 		2014-12-10 			inserita modifica								Tesser Paolo
% =================================================================================================
% 0.0.2			2014-12-11			inserita storia sull'introduzione				Luca S.
% =================================================================================================
% 0.0.3			2014-12-12			inserita storia sull'inizio cap. proc. prim.	Tesser Paolo
% =================================================================================================
% 0.0.4			2014-12-15			inserita storia sull'inizio cap. proc. org.		Tesser Paolo
% =================================================================================================
% 0.0.5			2014-12-26			inserita storia terminazione cap. proc. org.	Tesser Paolo
% =================================================================================================
% 1.0.1			2015-02-05			inserita histroy per piccole modifiche			Faccin Nicola
% =================================================================================================
% 1.0.2			2015-02-19			inserita history per aggiunta formazione		Faccin Nicola
% =================================================================================================
% 2.0.1			2015-04-02			inserita history per RP							Tesser Paolo
% =================================================================================================
%

\begin{center}
\begin{small}
	\textbf{\huge Diario Revisioni}
	\vspace{0.5cm}
	\begin{longtable}{p{6cm}|c|c|c}
		\label{tab:history}
		\textbf{Modifica} & \textbf{Autore \& Ruolo} & \textbf{Data} & \textbf{Versione} \\
		\hline

		% Fase di Progettazione di Detagglio - Requisiti opzionali - RQ





		% Fase di Progettazione di Detagglio - Requisiti desiderbaili





		% Fase di Progettazione di Detagglio - Requisiti obbligatori
		\emph{Approvazione documento} & 
			\begin{tabular}[c]{c c}
				Cusinato Giacomo \\
				\emph{Responsabile di Progetto} \\
		\end{tabular} & 2015-05-11 & v5.0.0 \\
		\hline
		\emph{Verifica del documento} &
			\begin{tabular}[c]{c c}
				Tesser Paolo \\
				\emph{Verificatore} \\
		\end{tabular} & 2015-05-11 & v4.1.0 \\
		\hline
		\emph{Ampliamento lista di controllo con errori tipici UML e codice in \ref{sec:lista_di_controllo}} &
			\begin{tabular}[c]{c c}
				Roetta Marco \\
				\emph{Amministratore di Progetto} \\
		\end{tabular} & 2015-05-08 & v4.0.1 \\
		\hline

		\emph{Approvazione documento} & 
			\begin{tabular}[c]{c c}
				Cusinato Giacomo \\
				\emph{Responsabile di Progetto} \\
		\end{tabular} & 2015-04-09 & v4.0.0 \\
		\hline
		\emph{Verifica del documento} &
			\begin{tabular}[c]{c c}
				Tesser Paolo \\
				\emph{Verificatore} \\
		\end{tabular} & 2015-04-08 & v3.1.0 \\
		\hline
		\emph{Inserimento strumento per la verifica Cloc in \ref{par:cloc}} &
			\begin{tabular}[c]{c c}
				Roetta Marco \\
				\emph{Amministratore di Progetto} \\
		\end{tabular} & 2015-04-07 & v3.0.5 \\
		\hline
		\emph{Aggiunti strumenti di test in Processi di Verifica: Wercker, Karma, Jasmine, Protractor, Testbed GAE} &
			\begin{tabular}[c]{c c}
				Roetta Marco \\
				\emph{Amministratore di Progetto} \\
		\end{tabular} & 2015-04-07 & v3.0.4 \\
		\hline
		\emph{Inserita appendice Configurazione File Python in \ref{sec:configurazione_file_python}} &
			\begin{tabular}[c]{c c}
				Roetta Marco \\
				\emph{Amministratore di Progetto} \\
		\end{tabular} & 2015-04-06 & v3.0.3 \\
		\hline
		\emph{Inserita appendice Configurazione File AngularJS in \ref{sec:configurazione_file_angularjs}} &
			\begin{tabular}[c]{c c}
				Roetta Marco \\
				\emph{Amministratore di Progetto} \\
		\end{tabular} & 2015-04-03 & v3.0.2 \\
		\hline
		\emph{Stesa sezione codifica e definizione prodotto in Processi di Sviluppo (Primari)} &
			\begin{tabular}[c]{c c}
				Roetta Marco \\
				\emph{Amministratore di Progetto} \\
		\end{tabular} & 2015-04-01 & v3.0.1 \\
		\hline


		% Fase di Progettazione Architetturale - RP
		\emph{Approvazione documento} & 
			\begin{tabular}[c]{c c}
				Santacatterina Luca \\
				\emph{Responsabile di Progetto} \\
		\end{tabular} & 2015-03-02 & v3.0.0 \\
		\hline
		\emph{Verifica del documento} &
			\begin{tabular}[c]{c c}
				Faccin Nicola \\
				\emph{Verificatore} \\
		\end{tabular} & 2015-02-26 & v2.1.0 \\
		\hline
		\emph{Inserito paragrafo sulle norme di codifica dei processi di sviluppo} &
			\begin{tabular}[c]{c c}
				Carnovalini Filippo \\
				\emph{Amministratore di Progetto} \\
		\end{tabular} & 2015-02-25 & v2.0.4 \\
		\hline
		\emph{Inserito paragrafo sul formalismo di classificazione dei test} &
			\begin{tabular}[c]{c c}
				Carnovalini Filippo \\
				\emph{Amministratore di Progetto} \\
		\end{tabular} & 2015-02-24 & v2.0.3 \\
		\hline
		\emph{Inserito attività riguardante la gestione della rotazione dei ruoli di progetto} &
			\begin{tabular}[c]{c c}
				Carnovalini Filippo \\
				\emph{Amministratore di Progetto} \\
		\end{tabular} & 2015-02-23 & v2.0.2 \\
		\hline
		\emph{Inseriti diagrammi per le procedure dei processi primari} &
			\begin{tabular}[c]{c c}
				Carnovalini Filippo \\
				\emph{Amministratore di Progetto} \\
		\end{tabular} & 2015-02-19 & v2.0.1 \\
		\hline


		% Fase di Analisi di Dettaglio - Interna
		\emph{Approvazione documento} & 
			\begin{tabular}[c]{c c}
				Ceccon Lorenzo \\
				\emph{Responsabile di Progetto} \\
		\end{tabular} & 2015-02-06 & v2.0.0 \\
		\hline
		\emph{Verifica del documento} &
			\begin{tabular}[c]{c c}
				Tesser Paolo \\
				\emph{Verificatore} \\
		\end{tabular} & 2015-02-05 & v1.1.0 \\
		\hline
		\emph{Redazione attività e norme di codifica e progettazione} &
			\begin{tabular}[c]{c c}
				Faccin Nicola\\
				\emph{Amministratore di Progetto} \\
		\end{tabular} & 2015-02-04 & v1.0.4 \\
		\hline
		\emph{Aggiunte strumento di chat Slack} &
			\begin{tabular}[c]{c c}
				Faccin Nicola \\
				\emph{Amministratore di Progetto} \\
		\end{tabular} & 2015-02-03 & v1.0.3 \\
		\hline
		\emph{Aggiunte attività di formazione nei processi organizzativi} &
			\begin{tabular}[c]{c c}
				Faccin Nicola \\
				\emph{Amministratore di Progetto} \\
		\end{tabular} & 2015-02-02 & v1.0.2 \\
		\hline
		\emph{Aggiornata struttura repository per le presentazioni} &
			\begin{tabular}[c]{c c}
				Faccin Nicola \\
				\emph{Amministratore di Progetto} \\
		\end{tabular} & 2015-02-01 & v1.0.1 \\
		\hline


		% Fase di Analisi dei Requisiti - RR
		\emph{Approvazione del documento} &
			\begin{tabular}[c]{c c}
				Cusinato Giacomo \\
				\emph{Responsabile di Progetto} \\
		\end{tabular} & 2014-12-23 & v1.0.0 \\
		\hline
		\emph{Verifica del documento} &
			\begin{tabular}[c]{c c}
				Roetta Marco \\
				\emph{Verificatore} \\
		\end{tabular} & 2014-12-20 & v0.2.0 \\
		\hline
		\emph{Correzione del documento} &
			\begin{tabular}[c]{c c}
				Santacatterina Luca \\
				\emph{Amministratore} \\
		\end{tabular} & 2014-12-19 & v0.1.1 \\
		\hline	
		\emph{Verifica del documento} &
			\begin{tabular}[c]{c c}
				Roetta Marco \\
				\emph{Verificatore} \\
		\end{tabular} & 2014-12-19 & v0.1.0 \\
		\hline		
		\emph{Completata stesura capitolo Processi Organizzativi} &
			\begin{tabular}[c]{c c}
				Tesser Paolo \\
				\emph{Amministratore} \\
		\end{tabular} & 2014-12-18 & v0.0.8 \\
		\hline
		\emph{Completata stesura capitolo Processi di Supporto} &
			\begin{tabular}[c]{c c}
				Santacatterina Luca \\
				\emph{Amministratore} \\
		\end{tabular} & 2014-12-16 & v0.0.7 \\
		\hline
		\emph{Iniziata stesura capitolo Processi Organizzativi} &
			\begin{tabular}[c]{c c}
				Tesser Paolo \\
				\emph{Amministratore} \\
		\end{tabular} & 2014-12-12 & v0.0.6 \\
		\hline
		\emph{Completata stesura capitolo Processi Primari} &
			\begin{tabular}[c]{c c}
				Tesser Paolo \\
				\emph{Amministratore} \\
		\end{tabular} & 2014-12-11 & v0.0.5 \\
		\hline
		\emph{Inizio stesura capitolo Processi Primari} &
			\begin{tabular}[c]{c c}
				Tesser Paolo \\
				\emph{Amministratore} \\
		\end{tabular} & 2014-12-09 & v0.0.4 \\
		\hline
		\emph{Iniziata stesura capitolo Processi di Supporto} &
			\begin{tabular}[c]{c c}
				Santacatterina Luca \\
				\emph{Amministratore} \\
		\end{tabular} & 2014-12-09 & v0.0.3 \\
		\hline
		\emph{Stesura capitolo Introduzione} &
			\begin{tabular}[c]{c c}
				Santacatterina Luca \\
				\emph{Amministratore} \\
		\end{tabular} & 2014-12-08 & v0.0.2 \\
		\hline
		\emph{Creazione struttura documento} &
			\begin{tabular}[c]{c c}
				Tesser Paolo \\
				\emph{Amministratore} \\
		\end{tabular} & 2014-12-08 & v0.0.1 \\
		\hline
		
	\end{longtable}

\end{small}
\end{center}
	\pagebreak
	\tableofcontents
	\pagebreak


	% PAGINE CONTENUTI
	\setcounter{page}{1}
	\pagestyle{styleDocPages}
	% =================================================================================================
% File:			content_files.tex
% Description:	Defiinisce i capitoli presenti nel documento
% Created:		2014-12-05
% Author:		Santacatterina Luca
% Email:		s88.luca@gmail.com
% =================================================================================================
% Modification History:
% Version		Modifier Date		Change											Author
% 0.0.1 		2014-12-05 			inserito capitolo di prova						Luca S.
% =================================================================================================
%

% DEFINIZIONE DEI CONTENUTI DEL DOCUMENTO

% =================================================================================================
% File:			nome_del_capitolo.tex
% Description:	Defiinisce i contenuti obbligatori della prima pagina
% Created:		2014/12/05
% Author:		Santacatterina Luca
% Email:		s88.luca@gmail.com
% =================================================================================================
% Modification History:
% Version		Modifier Date		Change											Author
% 0.0.1 		2014/12/05 			iniziata stesura documento di prova				Luca S.
% =================================================================================================
%

% CONTENUTO DEL CAPITOLO

\section{Capitolo di prova}

\subsection{Sotto capitolo di prova}
Lorem ipsum dolor sit amet, consectetur adipiscing elit, sed do eiusmod tempor incididunt ut labore et dolore magna aliqua. Ut enim ad minim veniam, quis nostrud exercitation ullamco laboris nisi ut aliquip ex ea commodo consequat. Duis aute irure dolor in reprehenderit in voluptate velit esse cillum dolore eu fugiat nulla pariatur. Excepteur sint occaecat cupidatat non proident, sunt in culpa qui officia deserunt mollit anim id est laborum.

\subsection{Sotto capitolo di prova}
Lorem ipsum dolor sit amet, consectetur adipiscing elit, sed do eiusmod tempor incididunt ut labore et dolore magna aliqua. Ut enim ad minim veniam, quis nostrud exercitation ullamco laboris nisi ut aliquip ex ea commodo consequat. Duis aute irure dolor in reprehenderit in voluptate velit esse cillum dolore eu fugiat nulla pariatur. Excepteur sint occaecat cupidatat non proident, sunt in culpa qui officia deserunt mollit anim id est laborum.

\subsection{Sotto capitolo di prova}
Lorem ipsum dolor sit amet, consectetur adipiscing elit, sed do eiusmod tempor incididunt ut labore et dolore magna aliqua. Ut enim ad minim veniam, quis nostrud exercitation ullamco laboris nisi ut aliquip ex ea commodo consequat. Duis aute irure dolor in reprehenderit in voluptate velit esse cillum dolore eu fugiat nulla pariatur. Excepteur sint occaecat cupidatat non proident, sunt in culpa qui officia deserunt mollit anim id est laborum.

\subsection{Sotto capitolo di prova}
Lorem ipsum dolor sit amet, consectetur adipiscing elit, sed do eiusmod tempor incididunt ut labore et dolore magna aliqua. Ut enim ad minim veniam, quis nostrud exercitation ullamco laboris nisi ut aliquip ex ea commodo consequat. Duis aute irure dolor in reprehenderit in voluptate velit esse cillum dolore eu fugiat nulla pariatur. Excepteur sint occaecat cupidatat non proident, sunt in culpa qui officia deserunt mollit anim id est laborum.

\end{document}